% --- LaTeX Essay Template - S. Venkatraman ---

% --- Set document class and font size ---

\documentclass[letterpaper, 11pt]{article}

% --- Package imports ---

% lipsum is just for generating placeholder text and can be removed
\usepackage{hyperref, lipsum} 

% --- Fonts (Set to Palatino or Times New Roman if desired) ---

% \usepackage{newpxtext, newpxmath}
% \usepackage{times}

% --- Page layout settiings ---

% Set page margins
\usepackage[left=1.2in, right=1.2in, top=.9in, bottom=.9in]{geometry}

% Anchor footnotes to the bottom of the page
\usepackage[bottom]{footmisc}
\usepackage{setspace}
\setstretch{1.15}
% Set line spacing
\renewcommand{\baselinestretch}{1.2}

% Use this for a one-off '*' footnote (e.g., a URL in a heading)
\newcommand{\starfootnote}[1]{%
\begingroup
\renewcommand{\thefootnote}{\fnsymbol{footnote}}%
\footnote{#1}%
\endgroup
\addtocounter{footnote}{-1}%
}

% Set spacing between paragraphs
\setlength{\parskip}{2mm}

% --- Page formatting settings ---



% --- Essay title and author ---
\title{\vspace{-1.5cm} Summary of \textit{Veridical Data Science}, Chapters 11 \footnote{\url{https://vdsbook.com/}}}

\date{\today}

% --- Main text ---
% --- Document starts here ---

\begin{document}
% Set link colors for labeled items and URLs (blue) 
\hypersetup{colorlinks=true, linkcolor=blue, urlcolor=blue}
\maketitle

\section*{Chapter 11: Binary Responses and Logistic Regression}

Chapter 11 introduces binary prediction problems and the tools used to evaluate them.

It explains how a continuous-response method like least squares can be transferred to classification.
It also suggests why LS can perform poorly for classification: it can produce predictions outside the [0,1] range, which is not interpretable as probabilities.

Thus, we need a different loss function and optimization method, which motivates logistic regression.
Rather than minimizing squared error, logistic regression does not require normality error assumptions
and minimizes the negative log-likelihood of the data, which is optimized with iterative numerical methods.

The chapter also discusses how to evaluate binary classifiers.
One metric is accuracy, which is the proportion of correct predictions. However, accuracy can be misleading in imbalanced datasets.
Using the confusion matrix, the chapter defines sensitivity (true positive rate) and specificity (true negative rate) and stresses that the appropriate tradeoff depends on problem context.
Another important metric is the area under the ROC curve (AUC), which measures the model's ability to rank positive instances higher than negative ones, regardless of the threshold.

Finally, the chapter applies PCS framework to logistic regression.
Both the logistic regression and LS model are easy to interpret, and we can check the stability of the coefficients by evaluating them across different subsets of data or different preprocessing steps.
To analyze the stability of the model, we can evaluate the metrics for different judgment calls, such as different thresholds for classification, different preprocessing steps, or different subsets of data.

% --- Document ends here ---
\end{document}
